\chapter{Nested TBA, Baxter Relations, and Separation Variables}
\label{ch:integrable-nested-tba-tq-separation-variables}

\ConventionsLorentzian

Diagonal TBA converts scalar Bethe--Yang equations into density and
pseudoenergy equations.  Non-diagonal scattering adds auxiliary roots.  The
thermodynamic limit then contains physical densities and magnonic auxiliary
densities.  This chapter records the precise structure needed later for
factorized relativistic models and for the planar gauge-theory integrability
volume.

\section{Nested Densities}

Start from logarithmic nested equations of the form
\[
  Lp_a(\theta_j)
  +
  \sum_{b,k}\delta_{ab}(\theta_j-\theta_k)
  +
  \sum_{\alpha,r}\chi_{a r}(\theta_j-u_\alpha^{(r)})
  =
  2\pi I_j ,
\]
together with auxiliary equations
\[
  \sum_{b,k}\chi_{r b}(u_\alpha^{(r)}-\theta_k)
  +
  \sum_{s,\beta}\psi_{rs}(u_\alpha^{(r)}-u_\beta^{(s)})
  =
  2\pi J_\alpha^{(r)} .
\]
The labels \(a,b\) enumerate physical particle species; \(r,s\) enumerate
auxiliary Dynkin nodes or nesting levels.  The functions \(\delta,\chi,\psi\)
are chosen branches of logarithms of the scalar and matrix-scattering
eigenvalue factors.

In the thermodynamic limit introduce physical densities
\(\rho_a,\rho_a^h\) and auxiliary densities
\(\sigma_r,\sigma_r^h\).  Differentiating the logarithmic equations gives
\begin{align}
  \rho_a+\rho_a^h
  &=
  \frac{1}{2\pi}\frac{\dd p_a}{\dd\theta}
  +\sum_b K_{ab}*\rho_b
  +\sum_r K_{a r}*\sigma_r ,
  \label{eq:nested-tba-physical-density}\\
  \sigma_r+\sigma_r^h
  &=
  \sum_a K_{r a}*\rho_a
  +\sum_s K_{rs}*\sigma_s .
  \label{eq:nested-tba-aux-density}
\end{align}
Here
\[
  (K*f)(x)=\int_\R K(x-y)f(y)\,\dd y,
\]
and each kernel is \(1/(2\pi)\) times the derivative of the corresponding
phase.  Auxiliary roots carry no direct one-particle energy in the examples
under discussion; they affect the free energy through the constraints.

\begin{proposition}[Nested TBA variational equations]
\label{prop:nested-tba-variational}
Assume the density equations
\eqref{eq:nested-tba-physical-density}--\eqref{eq:nested-tba-aux-density}
have a thermodynamic limit with positive total densities, and use the
fermionic occupancy entropy for each physical and auxiliary root species.
Define pseudoenergies
\[
  \epsilon_a=\log\frac{\rho_a^h}{\rho_a},
  \qquad
  \eta_r=\log\frac{\sigma_r^h}{\sigma_r}.
\]
Then stationarity of the constrained free energy at inverse temperature \(R\)
gives
\begin{align}
  \epsilon_a
  &=
  R E_a
  -\sum_b K_{ba}*\log(1+\ee^{-\epsilon_b})
  -\sum_r K_{r a}*\log(1+\ee^{-\eta_r}),
  \label{eq:nested-tba-epsilon}\\
  \eta_r
  &=
  -\sum_a K_{a r}*\log(1+\ee^{-\epsilon_a})
  -\sum_s K_{s r}*\log(1+\ee^{-\eta_s}).
  \label{eq:nested-tba-eta}
\end{align}
\end{proposition}

\begin{proof}
The proof is the same variational calculation as
Proposition~\ref{prop:tba-variational-derivation}, with a larger list of
densities.  Let the full list be \(n_A\), where \(A\) runs over both physical
and auxiliary labels, and let \(n_A^t=n_A+n_A^h\).  The entropy variation is
\[
  \delta s
  =
  \sum_A\int \left[
    \log\frac{n_A^h}{n_A}\,\delta n_A
    +
    \log\frac{n_A^t}{n_A^h}\,\delta n_A^t
  \right].
\]
The linear density constraints express every \(\delta n_A^t\) as a convolution
of kernels with variations of occupied densities.  Physical labels contribute
\(\delta(RE_a n_a)\) to the energy variation; auxiliary labels do not.  Since
the variations \(\delta\rho_a\) and \(\delta\sigma_r\) are independent, the
coefficient of each variation must vanish.  Substituting
\(\log(n_A^t/n_A^h)=\log(1+\ee^{-\epsilon_A})\) gives
\eqref{eq:nested-tba-epsilon}--\eqref{eq:nested-tba-eta}.
\end{proof}

\section{Constant A2 Magnonic Example}

The ultraviolet limit of many nested systems reduces to constant equations for
the plateau values \(Y_A=\ee^{\epsilon_A}\).  A minimal \(A_2\) magnonic
example has two auxiliary nodes and incidence matrix
\[
  I_{rs}=
  \begin{pmatrix}
    0&1\\
    1&0
  \end{pmatrix}.
\]
The constant Y-system is
\[
  Y_1^2=1+Y_2,\qquad
  Y_2^2=1+Y_1 .
\]
The symmetric solution has \(Y_1=Y_2=\phi_{\rm g}\), where
\[
  \phi_{\rm g}=\frac{1+\sqrt5}{2},\qquad
  \phi_{\rm g}^2=1+\phi_{\rm g}.
\]
This finite calculation is the algebraic core behind Rogers-dilogarithm
central-charge evaluations.  To turn it into a central-charge theorem for a
specific QFT one must also justify the scaling limit, the kernel reduction to
the constant Y-system, and the absence or subtraction of bulk-energy terms.

\section{Baxter TQ Relation}

For the \(XXX_{1/2}\) transfer eigenvalue
\eqref{eq:aba-xxx-transfer-eigenvalue}, define the Baxter polynomial
\[
  Q(u)=\prod_{j=1}^M(u-u_j).
\]
Then the eigenvalue equation can be rewritten as the finite-difference
relation
\begin{equation}
  \Lambda(u)Q(u)
  =
  \left(u+\frac{\ii}{2}\right)^L Q(u-\ii)
  +
  \left(u-\frac{\ii}{2}\right)^L Q(u+\ii).
  \label{eq:xxx-baxter-tq}
\end{equation}

\begin{proposition}[Bethe equations from the \(TQ\) relation]
\label{prop:tq-pole-bethe-equations}
Let \(Q(u)=\prod_j(u-u_j)\) have simple roots.  Equation
\eqref{eq:xxx-baxter-tq} has a polynomial \(\Lambda(u)\) if and only if the
roots \(u_j\) obey the \(XXX_{1/2}\) Bethe equations
\eqref{eq:aba-xxx-bethe-equations}.
\end{proposition}

\begin{proof}
If \(\Lambda\) is polynomial, evaluating \eqref{eq:xxx-baxter-tq} at
\(u=u_j\) gives
\[
  \left(u_j+\frac{\ii}{2}\right)^L Q(u_j-\ii)
  +
  \left(u_j-\frac{\ii}{2}\right)^L Q(u_j+\ii)=0.
\]
Since
\[
  \frac{Q(u_j+\ii)}{Q(u_j-\ii)}
  =
  -\prod_{\substack{k=1\\ k\neq j}}^M
  \frac{u_j-u_k+\ii}{u_j-u_k-\ii},
\]
the displayed equality is precisely
\eqref{eq:aba-xxx-bethe-equations}.  Conversely, if the Bethe equations hold,
the numerator on the right side of \eqref{eq:xxx-baxter-tq} vanishes at every
root of \(Q\), so division by \(Q\) produces a polynomial \(\Lambda\).
\end{proof}

For higher-rank symmetry, \(Q\)-functions are attached to nodes of the Dynkin
diagram, and the scalar \(TQ\) relation is replaced by a system of functional
relations.  The simplest form is a \(T\)-system or Hirota equation
\[
  T_{a,s}(u+\tfrac{\ii}{2})T_{a,s}(u-\tfrac{\ii}{2})
  =
  T_{a+1,s}(u)T_{a-1,s}(u)
  +
  T_{a,s+1}(u)T_{a,s-1}(u),
\]
with boundary conditions determined by the representation and by the physical
model.  The associated Y-functions are ratios of neighboring \(T\)'s.  These
functional equations organize both relativistic nested TBA systems and the
mirror-TBA/Y-system machinery used later in planar gauge theory.

\section{Q-Operators and Separation Variables}

A \(Q\)-operator is an operator-valued function \(Q(u)\) commuting with the
transfer matrix and satisfying an operator \(TQ\) relation whose eigenvalue
version is \eqref{eq:xxx-baxter-tq}.  Constructing \(Q(u)\) requires extra
auxiliary representations; the construction is model-dependent.  The
important logical point is that the Bethe roots are zeros of an eigenvalue of
an operator.  Their status comes from the spectral problem of the transfer
matrix.

Separation of variables uses a different complete set of commuting objects.
For the \(2\times2\) monodromy
\[
  T(u)=
  \begin{pmatrix}
    A(u)&B(u)\\
    C(u)&D(u)
  \end{pmatrix},
\]
one diagonalizes the commuting family \(B(u)\) in a representation where this
is possible.  If
\[
  B(u)=B_0\prod_{\alpha=1}^{L-1}(u-x_\alpha)
\]
on a separated eigenvector, the variables \(x_\alpha\) are separated
coordinates.  The transfer-matrix eigenvalue problem becomes a finite
difference equation in each \(x_\alpha\), whose one-variable form is the
Baxter equation.  The advantage is that a pseudovacuum need not exist.  The
cost is that the functional-analytic construction of the separated basis is
more delicate and depends on the representation.

\begin{openproblem}[Completeness beyond the pseudovacuum construction]
\label{op:bethe-completeness-sov}
Develop, for each integrable QFT family used in this monograph, a precise
finite-volume Hilbert-space theorem comparing algebraic Bethe ansatz,
\(Q\)-operator, and separation-of-variables bases.  The theorem should state
the representation space, boundary conditions, allowed singular roots, and the
topology in which the large-volume limit is taken.
\end{openproblem}

\section{Interface with Later Volumes}

The relativistic integrable-QFT use of nested TBA differs from planar
gauge-theory integrability.  In the present volume, physical particles have
relativistic dispersion \(E=m\cosh\theta\) and finite-size effects are
organized by the mirror channel of a two-dimensional QFT.  In the planar
\(\mathcal N=4\) SYM volume, magnons have a non-relativistic dispersion on a
spin chain, wrapping effects arise from trace cyclicity and finite operator
length, and the quantum spectral curve replaces the finite-rank Baxter
polynomial by a system of analytic \(Q\)-functions.  The algebraic bridge is
the same RTT, transfer-matrix, and \(TQ\) logic developed here; the physical
interpretation is different.
